% FAC-EN-U0048
% French source: fac_body.tex lines 1736--1764.
% Authority: PDF physical pp. 47--48; printed pp. 242--243 (no. 50 complete).
% The historical final paragraph is translated as printed; no modern result is
% silently substituted for Serre's 1955 statement.

\medskip
\phantomsection\label{fac:II:4:no:50}
\textbf{50. Projective modules and fiber spaces with vector-space fiber.}
\phantomsection\label{fac:II:4:no:50:def:projective-module}
Recall ([6], Chap.~I, Th.~2.2) that an $A$-module $M$ is called
\textit{projective} if it is a direct summand of a free $A$-module.

\phantomsection\label{fac:II:4:prop:4}
\textsc{Proposition 4.}
\textit{Let $M$ be a finitely generated $A$-module. For $M$ to be
projective, it is necessary and sufficient that the
$\mathcal{O}_x$-module $\mathcal{O}_x\otimes_A M$ be free for every
$x\in V$.}

If $M$ is projective, $\mathcal{O}_x\otimes_A M$ is
$\mathcal{O}_x$-projective, and hence $\mathcal{O}_x$-free since
$\mathcal{O}_x$ is a local ring (cf.~[6], Chap.~VIII, Th.~$6.1'$).

Conversely, if all the $\mathcal{O}_x\otimes_A M$ are free, then
\phantomsection\label{fac:II:4:prop:4:formula:projective-dimension}
\[
\dim(M)=\sup_{x\in V}\dim\bigl(\mathcal{O}_x\otimes_A M\bigr)=0
\qquad\text{(cf.~[6], Chap.~VII, Exer.~11),}
\]
which implies that $M$ is projective ([6], Chap.~VI, \S2).

Observe that if $\mathcal{F}$ is a coherent algebraic sheaf on $V$ and
$\mathcal{F}_x$ is isomorphic to $\mathcal{O}_x^p$, then $\mathcal{F}$ is
isomorphic to $\mathcal{O}^p$ over a neighbourhood of $x$. If this property
holds at every point $x\in V$, the sheaf $\mathcal{F}$ is therefore locally
isomorphic to a sheaf $\mathcal{O}^p$, with the integer $p$ remaining constant
on every connected component of $V$. Applying this to the sheaf
$\mathcal{A}(M)$ gives:

\phantomsection\label{fac:II:4:no:50:cor}
\textsc{Corollary.}
\textit{Let $\mathcal{F}$ be a coherent algebraic sheaf on a connected affine
variety $V$. The following three properties are equivalent:}

\phantomsection\label{fac:II:4:no:50:cor:i}
\noindent(i) \textit{$\Gamma(\mathcal{F})$ is a projective $A$-module.}

\phantomsection\label{fac:II:4:no:50:cor:ii}
\par\noindent(ii) \textit{$\mathcal{F}$ is locally isomorphic to a sheaf
$\mathcal{O}^p$.}

\phantomsection\label{fac:II:4:no:50:cor:iii}
\par\noindent(iii) \textit{$\mathcal{F}$ is isomorphic to the sheaf of germs of sections
of an algebraic fiber space with vector-space fiber and base $V$.}

\phantomsection\label{fac:II:4:no:50:correspondence}
In addition, the map
$E\longmapsto\Gamma\bigl(\mathcal{S}(E)\bigr)$ (where $E$ denotes a fiber
space with vector-space fiber) establishes a one-to-one correspondence
between the classes of fiber spaces and the classes of finitely generated
projective $A$-modules. Under this correspondence, a \textit{trivial} fiber
space corresponds to a \textit{free} module, and conversely.

\phantomsection\label{fac:II:4:no:50:historical-question}
Let us note that when $V=K^r$ (in which case
$A=K[X_1,\ldots,X_r]$), it is unknown whether there exist finitely generated
projective $A$-modules that are not free, or, equivalently, whether there
exist nontrivial algebraic fiber spaces with vector-space fiber and base
$K^r$.
